% !TeX spellcheck = ru_RU
% !TEX root = vkr.tex

\section{Экспериментальная часть}
\label{sec:experiment}

Экспериментальная проверка была направлена на работоспособность реорганизованной инфраструктуры в целом: корректность работы единого CLI, автоматизацию загрузки и подготовки матриц, единообразный запуск всех инструментов, а также сохранение результатов в унифицированном формате.

\subsection{Набор тестовых матриц}
\label{subsec:matrices}

Тестирование проводилось на матрицах из коллекции SuiteSparse Matrix Collection, используемых в предыдущих экспериментах по сравнению реализаций на основе взаимодействующих сетей. Характеристики матриц приведены в таблице~\ref{tab:matrices}.

\begin{table}[h]
    \centering
    \caption{Характеристики использованных матриц}
    \label{tab:matrices}
    \begin{tabular}{lrr}
        \hline
        Матрица & $|V|$ & $|E|$ \\
        \hline
        \texttt{usroads-48} & 126\,146 & 323\,900 \\
        \texttt{boneS01} & 127\,224 & 5\,516\,602 \\
        \texttt{coAuthorsCiteseer} & 227\,320 & 814\,134 \\
        \texttt{Andrews} & 60\,000 & 760\,154 \\
        \texttt{gr\_30\_30} & 900 & 7\,744 \\
        \texttt{Trefethen\_2000} & 2\,000 & 41\,906 \\
        \hline
    \end{tabular}
\end{table}

Для части матриц в экспериментах использовались как исходные, так и переформированные матрицы, полученные при помощи алгоритма обратного Катхилла--Макки (RCM), что позволяет оценить влияние переупорядочения на производительность рассматриваемых инструментов.

\subsection{Алгоритмы и конфигурация эксперимента}
\label{subsec:algorithms}

В экспериментах использовались три алгоритма: \emph{BFS}, \emph{SSSP} и \emph{Triangle Counting} (алгоритм Sandia). Набор матриц и соответствующих им алгоритмов задаётся в \texttt{matrices.yaml}: для каждой записи указывается имя матрицы, алгоритм, список инструментов для запуска, флаг \texttt{reorder}, а также ссылка для скачивания.

Конфигурация запуска бенчмарка задаётся через параметры командной строки: количество прогонов \texttt{run\_count}, список чисел потоков \texttt{thread\_counts} (либо значения по умолчанию, вычисляемые на основе числа доступных ядер процессора). Для каждого инструмента измерения проводятся отдельно, результаты сохраняются в индивидуальные файлы \texttt{results/raw/<timestamp>\_<tool>.json} в едином формате, описанном в разделе~\ref{subsec:common}.

\subsection{Воспроизведение эксперимента}
\label{subsec:reproducibility}

Воспроизведение эксперимента с использованием реорганизованной инфраструктуры осуществляется следующей последовательностью команд:

\begin{verbatim}
$ git submodule update --init --recursive
$ uv sync
$ uv run main.py download
$ uv run main.py prepare-matrices
$ uv run main.py benchmark inpla
$ uv run main.py benchmark fsharp
$ uv run main.py benchmark lagraph
$ uv run main.py benchmark networkx
\end{verbatim}

Единый интерфейс командной строки, централизованная конфигурация в \texttt{matrices.yaml}, автоматизация загрузки и подготовки матриц, единые утилиты запуска и единый формат результатов обеспечивают воспроизводимость экспериментов и упрощают добавление новых инструментов и наборов матриц.

Работоспособность инфраструктуры была проверена на всех перечисленных выше инструментах: для каждого из них настроена интеграция в соответствии с выбранным стилем, подготовка матриц выполнена автоматически, бенчмарки запущены через унифицированный CLI, результаты корректно сохранены в формате \texttt{Results}.
